% Gerado por src/python/thesis/build_tab_controlador.py -- nao editar a mao.
% Regenerar com: make thesis-tables
% Fonte: results/tables/controller_wtl_oos_20260326_221909.csv
% Fonte: data/processed/ppsn_controller_comparison_oos_20260326_221909.csv
% Fonte: data/processed/fla_response.csv
% Fonte: results/tables/dynamic_signal_lofo.csv
\begin{table}[!htb]
    \centering
    \small
    \caption{Controlador de aquecimento contra o SPEA2 e contra o IVF/SPEA2 sempre ativo, por classe de instância.}
    \label{tab:controlador_wtl}
    \begin{tabular*}{\textwidth}{@{\extracolsep{\fill}}llccc@{}}
        \toprule
        \multicolumn{3}{@{}c}{} & \multicolumn{2}{c@{}}{\textbf{V/E/D do controlador contra}} \\
        \cmidrule(l){4-5}
        \textbf{Métrica} & \textbf{Subconjunto} & \textbf{$n$} & \textbf{SPEA2 canônico} & \textbf{IVF/SPEA2 sempre ativo} \\
        \midrule
        \multirow{5}{*}{HV} & Todas & 51 & 34/14/3 & 2/40/9 \\
         & O operador ajuda & 41 & 33/8/0 & 0/32/9 \\
         & O operador não ajuda & 10 & 1/6/3 & 2/8/0 \\
         & \quad neutro & 6 & 1/5/0 & 0/6/0 \\
         & \quad prejudica & 4 & 0/1/3 & 2/2/0 \\
        \midrule
        \multirow{5}{*}{IGD} & Todas & 51 & 29/18/4 & 1/40/10 \\
         & O operador ajuda & 41 & 29/12/0 & 0/33/8 \\
         & O operador não ajuda & 10 & 0/6/4 & 1/7/2 \\
         & \quad neutro & 6 & 0/5/1 & 0/5/1 \\
         & \quad prejudica & 4 & 0/1/3 & 1/2/1 \\
        \bottomrule
    \end{tabular*}
    \par\smallskip
    \parbox{\textwidth}{\footnotesize
    V/E/D do controlador contra cada algoritmo, por subconjunto de instâncias; as linhas recuadas dividem as instâncias em que o operador não ajuda em neutras e prejudicadas (Seção~\ref{sec:ativacao}). Teste de postos sinalizados de Wilcoxon sobre 30 execuções pareadas por semente, com correção de Benjamini--Hochberg; um caso conta como vitória ou derrota quando o valor-$p$ ajustado fica abaixo de $0{,}05$ e o $A_{12}$ sai do intervalo $[0{,}44;\,0{,}56]$. Limiares de renovação aplicados: $0{,}216$ em DTLZ, MaF e ZDT, e $0{,}248$ em WFG.
    }
\end{table}
